% GENERATED FILE. Edit canonical lessons, then run npm run generate:books.
% canonical-source-hash: fnv1a64:fa8042482f535bbc
% canonical-lessons: ES-C17-futuro, ES-C17-comer-futuro, ES-C17-vivir-futuro

\chapter{The Future --- Saying What Will Happen}
\label{ch:spanish-futuro}
\hlchaptermodality{47}

\begin{chapteropening}
\textbf{By the end of this chapter:} \emph{I can say what I will do, in the singular, across all three verb families.}

The future endings ride on the whole infinitive, one family at a time.
\end{chapteropening}

\section[hablaré]{hablaré — keep the infinitive, add the future}

\label{lesson:ES-C17-futuro}



\begin{quote}
\textbf{Voy a hablar español} uses the near-future frame you already know. Spanish also has a one-word future: \textbf{Hablaré español} — “I will speak Spanish.” It can present a future prediction or intention without adding a new subject or object.
\end{quote}

\begin{grammarlens}[title={Grammar Lens: three familiar persons}]
Keep the whole infinitive \textbf{hablar}, then add one ending:

\noindent
\begin{tabularx}{\linewidth}{@{}>{\raggedright\arraybackslash}X>{\raggedright\arraybackslash}X@{}}
\toprule
\textbf{person} & \textbf{simple future} \\
\midrule
yo & hablar\textbf{é} \\
tú & hablar\textbf{ás} \\
él / ella / usted & hablar\textbf{á} \\
\bottomrule
\end{tabularx}

Say \textbf{hablaré · hablarás · hablará}. Each form has final stress, marked by an accent. Plural persons wait.
\end{grammarlens}

\begin{cousinweb}
Late Latin and early Romance used more than one way to speak about the future. One influential construction placed an infinitive beside present forms of Latin \textbf{habēre}, “to have.” Through repeated use and sound change, descendants of those pieces became the endings in forms such as \textbf{hablaré}.

That history explains why the infinitive remains visible. It does \textbf{not} mean that a modern learner should attach the full modern verb \textbf{haber}, and it does not require the claim that every older future vanished at one moment.
\end{cousinweb}

\subsection*{Your turn}
Say these aloud:

\begin{itemize}
\raggedright
  \item “hablaré, hablarás, hablará”
  \item “Voy a hablar español. Hablaré español.”
  \item “infinitive plus historical habēre material”
\end{itemize}

\begin{tcolorbox}[breakable,colback=teal!4,colframe=teal!35!black,title={Before you move on}]
Give the three learned forms. (\emph{Hablaré, hablarás, hablará.}) What stays visible? (The whole infinitive.) What older verb contributed to the endings? (Latin \textbf{habēre}.) Next: test the same endings on \textbf{comer}.
\end{tcolorbox}
\section[comeré]{comeré — the same endings keep -er}

\label{lesson:ES-C17-comer-futuro}



\begin{quote}
You built \textbf{hablaré} without removing \textbf{-ar}. Keep \textbf{comer} just as intact: \textbf{comeré} — “I will eat.”
\end{quote}

\begin{grammarlens}[title={Grammar Lens: the singular -er row}]
\noindent
\begin{tabularx}{\linewidth}{@{}>{\raggedright\arraybackslash}X>{\raggedright\arraybackslash}X@{}}
\toprule
\textbf{person} & \textbf{simple future} \\
\midrule
yo & comer\textbf{é} \\
tú & comer\textbf{ás} \\
él / ella / usted & comer\textbf{á} \\
\bottomrule
\end{tabularx}

Say \textbf{comeré · comerás · comerá}. Compare \textbf{hablaré / comeré}: the infinitives differ, but the future endings do not. Plural persons wait.
\end{grammarlens}

\begin{cousinweb}
The historical infinitive-plus-\textbf{habēre} construction helps explain the shared endings: the verb-family vowel remains inside \textbf{hablar} or \textbf{comer}, rather than selecting a different ending set.
\end{cousinweb}

\subsection*{Your turn}
Say these aloud:

\begin{itemize}
\raggedright
  \item “comeré, comerás, comerá”
  \item “Hablaré. Comeré.”
  \item “Beberé café.”
\end{itemize}

\begin{tcolorbox}[breakable,colback=teal!4,colframe=teal!35!black,title={Before you move on}]
Give the three learned forms. (\emph{Comeré, comerás, comerá.}) Did the \textbf{-er} disappear? (No.) Which endings match \textbf{hablar}? (\textbf{-é, -ás, -á}.) Next: let \textbf{vivir} test the pattern once more.
\end{tcolorbox}
\section[viviré]{viviré — one future row for three verb families}

\label{lesson:ES-C17-vivir-futuro}



\begin{quote}
\textbf{Hablaré} kept \textbf{-ar}. \textbf{Comeré} kept \textbf{-er}. Predict the yo future of \textbf{vivir} before reading on. It is \textbf{viviré}.
\end{quote}

\begin{grammarlens}[title={Grammar Lens: the singular -ir row}]
\noindent
\begin{tabularx}{\linewidth}{@{}>{\raggedright\arraybackslash}X>{\raggedright\arraybackslash}X@{}}
\toprule
\textbf{person} & \textbf{simple future} \\
\midrule
yo & vivir\textbf{é} \\
tú & vivir\textbf{ás} \\
él / ella / usted & vivir\textbf{á} \\
\bottomrule
\end{tabularx}

Say \textbf{viviré · vivirás · vivirá}. Now the bounded pattern is complete: \textbf{hablaré, comeré, viviré}. The same three endings attach to all three whole infinitives. Plural persons wait.
\end{grammarlens}

\begin{grammarlens}[title={Grammar Lens: meaning with known words}]
\textbf{Viviré en Madrid} says “I will live in Madrid.” Every lexical piece was available before this lesson; only the future form is new.
\end{grammarlens}

\subsection*{Your turn}
Say these aloud:

\begin{itemize}
\raggedright
  \item “viviré, vivirás, vivirá”
  \item “hablaré, comeré, viviré”
  \item “Viviré en Madrid.”
\end{itemize}

\begin{tcolorbox}[breakable,colback=teal!4,colframe=teal!35!black,title={Before you move on}]
Give the three \textbf{vivir} forms. (\emph{Viviré, vivirás, vivirá.}) Which verb families share the learned endings? (All three: \textbf{-ar, -er, -ir}.) Next: reuse the whole-infinitive frame for “would.”
\end{tcolorbox}
